\chapter{Week 8: Macros}
\label{ch:week8}

Lecture 8: October 27, 2026. Recitation 8: October 30, 2026.

Assigned reading and resources (course page):
\begin{itemize}
  \item \cite{PLAI} Chapter 13 (Desugaring as a Language Feature).
  \item \cite{RacketGuide} Chapter 16 (Macros), and the further resources it links to.
  \item \cite{ClingerWand}, optional: the history of how hygienic macro systems were developed (\Cref{sec:hygiene}).
  \item Lecture code: \code{lecture08.rkt}.
\end{itemize}

Further reading:
\begin{itemize}
  \item \cite{FearOfMacros}, a tutorial on \code{syntax-case}, \code{with-syntax}, and \code{datum->syntax}, recommended by the Racket Guide; it covers deliberately unhygienic macros like \code{aif} (\Cref{sec:aif}).
  \item \cite{PLAI} Chapter 4 (A First Taste of Desugaring), the desugaring inside an interpreter that macros hand to the programmer.
  \item \cite{TAPL} Section 5.2 (Programming in the Lambda-Calculus): Church booleans and numerals, and recursion, as in Recitation~8 (\Cref{sec:rec8}); optional.
\end{itemize}

Some constructs cannot be functions. A function call evaluates every argument first, once, so a function cannot skip an argument (\code{or}, \code{if}), repeat one (\code{while}), look at an argument's text (\code{assert}), or assign to the caller's variables (\code{swap}). A \emph{macro} can: it is a function from syntax to syntax that Racket runs before evaluation, rewriting each use into other code. This is the desugaring our interpreters have done all along (\code{let} into \code{lambda}, \code{rec} into boxes), now available to the programmer. Hygiene can also be broken on purpose, one name at a time. The recitation turns away from macros to the lambda calculus, a language with nothing but functions, and encodes booleans, numbers, and recursion in it. This chapter covers:
\begin{itemize}
  \item \Cref{sec:orfun}: why \code{or} cannot be a function under call-by-value;
  \item \Cref{sec:myor}: \code{define-syntax}, \code{syntax-case}, syntax objects and \code{\#'} templates, with \code{my-or}, \code{my-if}, \code{define-syntax-rule}, and \code{unless};
  \item \Cref{sec:mywhile}: a recursive macro that never stops expanding, and a \code{while} loop that expands into a recursive function instead;
  \item \Cref{sec:myassert}: a macro that uses its argument both as code to run and as data to print;
  \item \Cref{sec:multior}: ellipsis patterns, recursion that terminates by shrinking the arguments, and binding an argument once to avoid evaluating it twice;
  \item \Cref{sec:hygiene}: variable capture, hygienic expansion, and hygiene as lexical scope for macros;
  \item \Cref{sec:literals}: keywords in patterns with the literal list;
  \item \Cref{sec:mylet}: \code{let}, \code{let*}, and multi-binding \code{let} as macros into \code{lambda};
  \item \Cref{sec:swap}: \code{swap} as a function, with boxes, and as a macro that mutates the caller's variables;
  \item \Cref{sec:aif}: an anaphoric \code{aif} that deliberately captures \code{it} with \code{datum->syntax} (not covered in lecture);
  \item \Cref{sec:rec8}: Recitation~8, the lambda calculus: substitution and alpha-equivalence, normal-order evaluation, Church booleans and numerals, and \code{Y} under normal order.
\end{itemize}

\section{Why \texttt{or} Is Not a Function}
\label{sec:orfun}

Racket's \code{or} stops at the first argument that is true. In
\begin{lstlisting}
(or #t (print "Hello World!"))
\end{lstlisting}
the first argument is already \code{\#t}, so the result is \code{\#t} and the \code{print} never runs: nothing is printed.

The obvious attempt to write our own \code{or} is a function:
\begin{lstlisting}
(define (my-or1 a b)
  (if a #t b))
\end{lstlisting}
Its body looks right: if \code{a} is true, answer \code{\#t} without looking at \code{b}. But the call prints:
\begin{lstlisting}[style=repl]
> (my-or1 #t (print "hello world"))
"hello world"#t
\end{lstlisting}
\code{print} writes the string with its quotes and no newline, and the REPL then prints the result \code{\#t} right after it. Writing the body out by hand, with the argument expression in place of \code{b}, does not print:
\begin{lstlisting}[style=repl]
> (if #t #t (print "hello world"))
#t
\end{lstlisting}
So the function and its inlined body behave differently. Racket evaluates every argument of a function call before the body runs (eager, call-by-value evaluation, as in our interpreters since \Cref{sec:closures}). By the time the \code{if} inside \code{my-or1} looks at \code{a}, the side effect in \code{b} has already happened. The body cannot undo that.

\begin{keypoint}
A function receives values, not expressions. \code{or} has to decide \emph{whether} to evaluate its second argument, so it cannot be a function under call-by-value. It has to be a rewrite of the program text, done before evaluation: \code{(or a b)} becomes something like \code{(if a \#t b)}, where \code{b} sits in a branch of the \code{if} and is evaluated only if that branch is taken.
\end{keypoint}

The extra evaluation shows up through effects: printing, mutation, errors, and non-termination. For example, \code{(my-or1 \#t (car '()))} raises an error where \code{(or \#t (car '()))} returns \code{\#t}.

\begin{remark}
\code{my-or1} also differs from \code{or} in the value it returns: Racket's \code{or} returns the first true value itself, so \code{(or 5 \#f)} is \code{5}, while \code{(my-or1 5 \#f)} is \code{\#t}.
\end{remark}

\section{A First Macro: \texttt{my-or}}
\label{sec:myor}

\Cref{sec:orfun} showed that \code{or} must rewrite the program before it runs. Racket lets us write such a rewrite ourselves, as a \emph{macro}:
\begin{lstlisting}
(define-syntax (my-or stx)
  (syntax-case stx ()
    [(my-or a b)
     #'(if a #t b)]))
\end{lstlisting}
Piece by piece:
\begin{itemize}
  \item \code{define-syntax} binds \code{my-or} as a macro instead of a function. Whenever the expander meets a form \code{(my-or \dots)}, it calls the macro's \emph{transformer} on that form, before the program runs.
  \item The transformer is a function. Its argument \code{stx} is a \emph{syntax object}: the program text of the whole use, such as \code{(my-or \#t (print "hello world"))}, together with information like source location and scope. It returns another syntax object, which replaces the use.
  \item \code{syntax-case} takes \code{stx} apart by pattern matching. The \code{()} is the list of literals, keywords that must match exactly; there are none here. The pattern \code{(my-or a b)} matches a use with exactly two subforms and binds the \emph{pattern variables} \code{a} and \code{b} to them, unevaluated. The first position, \code{my-or}, is a pattern variable too: it matches the macro's name as written at the use, and is usually ignored (it is often written \code{\_}).
  \item \code{\#'(if a \#t b)} is short for \code{(syntax (if a \#t b))}: a \emph{template}. It builds a syntax object in which every pattern variable is replaced by the piece of syntax it matched. \code{\#'} is to syntax objects what \code{'} is to lists.
\end{itemize}

So the use
\begin{lstlisting}
(my-or #t (print "hello world"))
\end{lstlisting}
is rewritten into
\begin{lstlisting}
(if #t #t (print "hello world"))
\end{lstlisting}
which is exactly the expression from \Cref{sec:orfun} that returns \code{\#t} without printing. The argument \code{(print "hello world")} is moved into a branch of the \code{if} as program text, never as a value, so it is evaluated only if that branch is taken.

\begin{keypoint}
A function maps values to a value at run time. A macro maps syntax to syntax at expansion time, before evaluation starts. This is the same desugaring we did inside the interpreter, \code{let} into an application of a \code{lambda} (\Cref{sec:substrev}) and \code{rec} into boxes (\Cref{sec:recboxes}), now done by the language for the programmer.
\end{keypoint}

\subsection{\texttt{my-if}}

The same recipe gives a conditional of our own, written in terms of \code{cond}:
\begin{lstlisting}
(define-syntax (my-if stx)
  (syntax-case stx ()
    [(my-if a b c)
     #'(cond [a b] [else c])]))
\end{lstlisting}
\code{(my-if e1 e2 e3)} expands to \code{(cond [e1 e2] [else e3])}, so only one of \code{e2} and \code{e3} is evaluated. As a function, \code{(define (my-if a b c) (cond [a b] [else c]))} would evaluate both branches before choosing one. Every effect in the branch not taken would happen anyway, and a recursive function whose base case is guarded by such an \code{if}, like
\begin{lstlisting}
(define (fact n)
  (my-if (= n 0) 1 (* n (fact (- n 1)))))
\end{lstlisting}
would never terminate: the recursive call in the third argument is evaluated even when \code{n} is 0. With the macro, \code{fact} works, because the recursive call ends up in the \code{else} branch of the \code{cond}.

\begin{remark}
In the template, \code{else} is not a pattern variable: it is not in the pattern \code{(my-if a b c)}, so it is copied into the output as is, and \code{cond} sees it as its own \code{else} keyword. Only \code{a}, \code{b}, and \code{c} are replaced.
\end{remark}

\subsection{\texttt{define-syntax-rule}}

When a macro has a single pattern, \code{define-syntax-rule} says the same thing more briefly:
\begin{lstlisting}
(define-syntax-rule (my-or2 a b)
  (if a #t b))
\end{lstlisting}
The first part is the pattern and the second is the template. There is no \code{stx}, no \code{syntax-case}, and no \code{\#'}: the template is implicitly a \code{syntax} template. \code{my-or2} behaves exactly like \code{my-or}. The longer \code{syntax-case} form is needed when a macro has several patterns, or when its transformer must run code of its own before producing the output.

\subsection{\texttt{unless}}

\code{(unless b e)} evaluates \code{e} only when \code{b} is false. It is \code{if} with the branches swapped and nothing in the true branch:
\begin{lstlisting}
(define-syntax-rule (unless b e)
  (if b (void) e))
\end{lstlisting}
\code{(void)} returns Racket's ``no useful value'' value, which the REPL does not print. As with \code{or} and \code{if}, \code{unless} must be a macro: as a function, \code{e} would be evaluated before the test, so \code{(unless \#t (print "hi"))} would print.

\begin{remark}
Racket already has an \code{unless}; a definition at the top of a module shadows the imported one. Racket's version also allows several body expressions, \code{(unless b e1 e2 \dots)}, which needs an ellipsis pattern, \code{(unless b e ...)}.
\end{remark}

\section{A Loop: \texttt{my-while}}
\label{sec:mywhile}

\code{(my-while test body)} should evaluate \code{body} again and again as long as \code{test} is true. Like \code{or}, it cannot be a function: \code{test} and \code{body} must be evaluated many times, and a function receives each argument once, already evaluated.

\subsection{A Macro That Never Stops Expanding}

The first attempt mirrors the definition of a while loop: if the test holds, run the body and then the whole loop again.
\begin{lstlisting}
(define-syntax-rule (my-while test body)
  (cond [test (begin body (my-while test body))] [else (void)]))
\end{lstlisting}
The definition itself is accepted. The problem appears at the first use. The expander replaces \code{(my-while test body)} by the template, and the template contains \code{(my-while test body)} again, which the expander must also replace, and so on. Expansion happens before the program runs, so \code{test} has no value yet and nothing can stop the unfolding: expansion never terminates, and the program never starts.

\begin{keypoint}
A macro that uses itself in its template must make the recursive use \emph{smaller} (fewer arguments, a shorter list), so that expansion reaches a base case. Repetition that depends on run-time values, such as a loop condition, cannot be done by expansion. It has to be done by run-time code that the macro generates.
\end{keypoint}

\subsection{Generating a Recursive Function}

The fix moves the recursion from expansion time to run time. The macro expands, once, into a recursive function and a call to it:
\begin{lstlisting}
(define-syntax-rule (my-while test body)
  (begin (define (loop)
           (if test (begin body (loop)) (void)))
         (loop)))
\end{lstlisting}
Here \code{(loop)} inside the template is an ordinary function call, not a macro use, so the expander leaves it alone. At run time, each call of \code{loop} evaluates \code{test} afresh; if it is true, it runs \code{body} and calls \code{loop} again, and if it is false, it returns \code{(void)}. The \code{test} and \code{body} expressions are pasted into the body of \code{loop}, which is exactly what makes them run once per iteration.

Using it:
\begin{lstlisting}
(define b (box 0))

(my-while (< (unbox b) 3)
          (begin (print "Hello!")
                 (set-box! b (+ 1 (unbox b)))))
\end{lstlisting}
This prints \code{"Hello!"} three times, and \code{b} ends up holding 3.

\begin{remark}
The template defines a function named \code{loop}, and the user's \code{test} or \code{body} might mention a \code{loop} of their own. By hygiene (\Cref{sec:hygiene}), the \code{loop} written in the template and a \code{loop} written by the user are different variables, so the macro's helper cannot capture the user's name.
\end{remark}

\section{Macros See the Code: \texttt{my-assert}}
\label{sec:myassert}

The macros so far needed control over \emph{when} an argument is evaluated. A macro can also use the argument's \emph{text}:
\begin{lstlisting}
(define-syntax-rule (my-assert e)
  (unless e (printf "assertion failed: ~s\n" 'e)))
\end{lstlisting}
In the template, \code{e} appears twice. The first \code{e} is evaluated as the condition. The second is inside a quote: \code{'e} is \code{(quote e)}, and since \code{e} is a pattern variable it is replaced by the user's expression first, so the quote produces that expression as data. An assertion that holds, such as
\begin{lstlisting}
(my-assert (= (+ 1 1) 2))
\end{lstlisting}
prints nothing: the condition is true, so \code{unless} returns \code{(void)}. A failing one,
\begin{lstlisting}
(my-assert (= (+ 1 1) 3))
\end{lstlisting}
expands to
\begin{lstlisting}
(unless (= (+ 1 1) 3) (printf "assertion failed: ~s\n" '(= (+ 1 1) 3)))
\end{lstlisting}
and prints \code{assertion failed: (= (+ 1 1) 3)}: the expression as written, not its value \code{\#f} and not a simplified form. The \code{\textasciitilde s} directive of \code{printf} prints a value the way \code{write} does.

\begin{keypoint}
A function \code{(define (my-assert v) \dots)} would receive only \code{\#f} and could not report which assertion failed. A macro receives the expression itself, so it can both evaluate it and print it.
\end{keypoint}

\begin{remark}
In this file, \code{unless} is the macro defined earlier, which takes exactly one body expression. The template uses it with one, so it works with either version.
\end{remark}

\section{Any Number of Arguments: \texttt{multi-or}}
\label{sec:multior}

Racket's \code{or} takes any number of arguments. A macro accepts that with an \emph{ellipsis} pattern, and handles it by recursion on the number of arguments:
\begin{lstlisting}
(define-syntax (multi-or stx)
  (syntax-case stx ()
    [(multi-or e0 e1 ...)
     #'(let ([b e0])
         (if b b (multi-or e1 ...)))]
    [(multi-or)
     #'#f]))
\end{lstlisting}

\subsection{Ellipses}

In a pattern, \code{e1 ...} matches zero or more subforms, and binds \code{e1} to the whole sequence. So \code{(multi-or e0 e1 ...)} matches any use with at least one argument: \code{e0} is the first, \code{e1} the rest. In a template, \code{e1 ...} puts that sequence back, in order. The second clause, \code{(multi-or)}, matches the use with no arguments, and its template \code{\#'\#f} is the answer of an empty \code{or}. \code{syntax-case} tries the clauses in order, and the two patterns do not overlap.

\subsection{Recursion That Terminates}

The template uses \code{multi-or} again, as \code{my-while} did in \Cref{sec:mywhile}, but here the recursive use has one argument fewer. Every expansion step removes one argument, so expansion reaches \code{(multi-or)} and stops. For example:
\begin{lstlisting}
(multi-or #f #f 3)
---> (let ([b #f]) (if b b (multi-or #f 3)))
---> (let ([b #f]) (if b b (let ([b #f]) (if b b (multi-or 3)))))
---> ... (let ([b 3]) (if b b (multi-or))) ...
---> ... (let ([b 3]) (if b b #f)) ...
\end{lstlisting}
The result is \code{3}. Like Racket's \code{or}, and unlike \code{my-or} of \Cref{sec:myor}, \code{multi-or} returns the first true value itself, not \code{\#t}.

\subsection{Why the \texttt{let}}

The lecture code also defines the version without the \code{let}, and calls it \code{multi-or-bad}:
\begin{lstlisting}
(define-syntax (multi-or-bad stx)
  (syntax-case stx ()
    [(multi-or-bad e0 e1 ...)
     #'(if e0 e0 (multi-or-bad e1 ...))]
    [(multi-or-bad)
     #'#f]))
\end{lstlisting}
Its template copies \code{e0} twice, so \code{e0} is evaluated twice when it is true. With effects that is visible: \code{(multi-or-bad (begin (print "hi") 1))} prints \code{"hi"} twice, once for the test and once for the result. \code{multi-or} binds \code{e0} once to \code{b} and then tests and returns \code{b}, so it evaluates \code{e0} exactly once.

The same issue arises in any macro that uses an argument twice. A \code{square} macro written as \code{(* x x)} would evaluate its argument twice, so it binds it first:
\begin{lstlisting}
(define-syntax-rule (square x)
  (let ([v x])
    (* v v)))
\end{lstlisting}

\begin{keypoint}
A pattern variable in a template stands for a piece of code, not a value. Using it twice in the template copies the code twice, and the code then runs twice. When a macro needs an argument's value more than once, it should bind the argument to a variable once.
\end{keypoint}

\subsection{A Name Introduced by the Macro}

The \code{let} introduces a name, \code{b}, that the user never wrote. What happens when the user's own arguments mention a \code{b} is the subject of \Cref{sec:hygiene}.

\section{Hygiene}
\label{sec:hygiene}

\subsection{The Problem: Capture}

In lecture, \code{multi-or} of \Cref{sec:multior} was used under a binding of the user's own named \code{b}:
\begin{lstlisting}
(let ([b #t])
  (multi-or #f b))
\end{lstlisting}
If expansion just pasted text, as the C preprocessor does (\Cref{sec:rec5}), the result would be the expansion written out in a comment in the lecture code:
\begin{lstlisting}
(let ([b #t])
  (let ([b #f])                 ; b written in the template
    (if b b (multi-or b))))     ; b written by the user
\end{lstlisting}
The user's \code{b}, which referred to the outer \code{\#t}, now sits inside the macro's \code{(let ([b \#f]) \dots)} and refers to \code{\#f}. The program answers \code{\#f} although one argument of the \code{or} is true. The user's variable has been \emph{captured} by a binding the user cannot see.

\subsection{Hygienic Expansion}

\begin{definition}[Hygienic macro system]
A macro system is \emph{hygienic} if expansion preserves lexical scope: every name in the expanded program refers to the binding that was visible where the name was \emph{written}. A name written in a macro's template is resolved where the macro is defined; a name written by the user in an argument is resolved where the macro is used.
\end{definition}

Racket's expander is hygienic. It records, for every identifier, whether it came from the template or from the use, so the two \code{b}s above are different variables even though they print the same. The effect is as if the template's \code{b} had been renamed to a fresh name before pasting,
\begin{lstlisting}
(let ([b #t])
  (let ([b1 #f])
    (if b1 b1 (multi-or b))))
\end{lstlisting}
and the program answers \code{\#t}, which is what DrRacket printed in lecture.

\subsection{Two Directions of Capture}

Hygiene rules out capture in both directions.
\begin{enumerate}
  \item \emph{The macro captures the user's name.} A binding introduced by the template, like \code{b} above, must not capture a variable the user wrote in an argument.
  \item \emph{The user captures the macro's name.} A free name in the template, like \code{if} or \code{let}, must not be captured by a binding at the use site. In
\begin{lstlisting}
(let ([if 42])
  (multi-or #f 1))
\end{lstlisting}
  the user has rebound \code{if} locally, but the \code{if} in the template still means Racket's \code{if}, the one visible where \code{multi-or} was defined. The answer is \code{1}.
\end{enumerate}
The \code{loop} helper of \code{my-while} (\Cref{sec:mywhile}) is protected by the first rule: a user variable named \code{loop} in \code{test} or \code{body} is not the macro's \code{loop}.

\subsection{Connections}

\begin{keypoint}
Hygiene is lexical scope for macros. A text macro resolves every name after pasting, at the use site, which is the macro analogue of dynamic scope (\Cref{sec:dynscope}), with the same symptom: renaming a variable changes what the program does. A hygienic expander avoids capture the way capture-avoiding substitution did in \Cref{sec:capavoid}, by giving the introduced names fresh identities.
\end{keypoint}

\begin{remark}
DrRacket's \emph{Macro Stepper} shows each expansion step. On the example above, it displays the template's \code{b} and the user's \code{b} in different colors, which is a quick way to see which name refers to which binding.
\end{remark}

\section{Keywords in Patterns: \texttt{my-if2}}
\label{sec:literals}

Every \code{syntax-case} so far had an empty list \code{()} after \code{stx}. That list names \emph{literals}: identifiers that a pattern matches exactly, instead of binding them as pattern variables. They let a macro have keywords of its own:
\begin{lstlisting}
(define-syntax (my-if2 stx)
  (syntax-case stx (then else)
    [(my-if2 e1 then e2 else e3)
     #'(if e1 e2 e3)]))
\end{lstlisting}
In the pattern, \code{e1}, \code{e2}, and \code{e3} are pattern variables, but \code{then} and \code{else} are listed as literals, so the use must contain those words in those positions:
\begin{lstlisting}
(my-if2 (> 2 1) then "yes" else "no")    ; expands to (if (> 2 1) "yes" "no")
\end{lstlisting}
A use without the keywords, such as \code{(my-if2 (> 2 1) "yes" "no")}, matches no clause, and the expander reports a syntax error for \code{my-if2} before the program runs.

\begin{pitfall}[Forgetting the literal list]
Without \code{(then else)}, the pattern \code{(my-if2 e1 then e2 else e3)} would treat \code{then} and \code{else} as two more pattern variables. Every use with five subforms would match, whatever words stood in those positions, and \code{(my-if2 c 1 2 3 4)} would expand to \code{(if c 2 4)} without complaint.
\end{pitfall}

\begin{remark}
A literal is matched as an identifier, by binding, not only by spelling: the \code{else} in the use must refer to the same \code{else} as the one in the macro's literal list. This is hygiene again (\Cref{sec:hygiene}), applied to keywords. \code{cond}'s \code{else} works the same way.
\end{remark}

\section{\texttt{let} as a Macro}
\label{sec:mylet}

In \Cref{sec:substrev} we observed that \code{let} is an application of a \code{lambda} in disguise, and our interpreters used that equation to desugar \code{letC}. A macro states the same equation in Racket itself:
\begin{lstlisting}
(define-syntax (my-let1 stx)
  (syntax-case stx ()
    [(my-let1 ([var val]) body)
     #'((lambda (var) body) val)]))
\end{lstlisting}
The pattern matches a \code{let} with exactly one binding, written with Racket's usual double parentheses. For example,
\begin{lstlisting}
(my-let1 ([x 3]) (+ x x))
---> ((lambda (x) (+ x x)) 3)
\end{lstlisting}
which evaluates to 6. Square brackets and parentheses are interchangeable in patterns, as everywhere in Racket; the brackets only make the binding easier to read.

The pattern variable \code{var} appears in the template in a \emph{binding} position, as the parameter of the \code{lambda}. So the user's \code{x} is bound by a \code{lambda} the macro wrote. This is fine under hygiene: the user wrote both the binding \code{x} and its uses in \code{body}, so they agree, and they are both resolved at the use site.

\subsection{\texttt{my-let*}: Recursion over the Bindings}

\code{let*} binds its variables one at a time, so each right-hand side can use the variables before it. Desugared into \code{lambda}s, each binding becomes one application, nested inside the previous one:
\begin{lstlisting}
(my-let* ([x1 e1] [x2 e2]) e3)
---> ((lambda (x1) ((lambda (x2) e3) e2)) e1)
\end{lstlisting}
Since \code{e2} sits inside \code{(lambda (x1) \dots)}, it is in the scope of \code{x1}. The macro produces this nesting by recursion on the list of bindings, the same shape as \code{multi-or} (\Cref{sec:multior}):
\begin{lstlisting}
(define-syntax (my-let* stx)
  (syntax-case stx ()
    [(my-let* ([var val]) body)
     #'((lambda (var) body) val)]
    [(my-let* ([var val] . rest) body)
     #'((lambda (var) (my-let* rest body)) val)]))
\end{lstlisting}
The first clause is the base case, a single binding, which is \code{my-let1}. The second peels off the first binding. Its pattern \code{([var val] . rest)} is a dotted pair, as in \Cref{ch:week1}: \code{rest} matches the \emph{list} of the remaining bindings, so the template can pass it on unchanged, \code{(my-let* rest body)}, to a recursive use with one binding fewer. Expansion therefore terminates. An equivalent pattern is \code{([var val] rest ...)}, where \code{rest} matches the remaining bindings as a sequence and the template must wrap them back into a list, \code{(my-let* (rest ...) body)}. The order of the clauses matters: a single binding also matches the second pattern, with \code{rest} empty, and would expand to \code{(my-let* () body)}, which no clause matches. Listing the one-binding clause first catches that case before it reaches the second.

For example, \code{(my-let* ([x 1] [y (+ x 1)]) (* x y))} expands to \code{((lambda (x) ((lambda (y) (* x y)) (+ x 1))) 1)} and evaluates to 2. With \code{let} instead of \code{let*}, the \code{x} in \code{(+ x 1)} would be unbound, because all right-hand sides of a \code{let} are evaluated outside the new bindings.

\subsection{\texttt{my-let}: Many Bindings at Once}

A \code{let} with several bindings is a single \code{lambda} with several parameters, applied to all the right-hand sides at once:
\begin{lstlisting}
(let ([x1 e1] [x2 e2] [x3 e3]) body)
---> ((lambda (x1 x2 x3) body) e1 e2 e3)
\end{lstlisting}
No recursion is needed; an ellipsis does it in one step:
\begin{lstlisting}
(define-syntax (my-let stx)
  (syntax-case stx ()
    [(my-let ([var val] ...) body)
     #'((lambda (var ...) body) val ...)]))
\end{lstlisting}
In the pattern, the ellipsis follows the sub-pattern \code{[var val]}, so it matches a sequence of pairs, and \code{var} and \code{val} each stand for a sequence: all the names, and all the right-hand sides, in order. The template uses the two sequences separately, \code{var ...} as the parameter list and \code{val ...} as the arguments. The ellipsis thus ``unzips'' the list of pairs into two lists.

\begin{keypoint}
\code{let} and \code{let*} desugar differently. In \code{my-let}, every \code{val} is an argument of the one application, outside the \code{lambda}, so no right-hand side can see any of the new variables. In \code{my-let*}, each \code{val} is inside the \code{lambda}s of the earlier bindings, so it can see them.
\end{keypoint}

\section{\texttt{swap}: A Macro That Mutates Variables}
\label{sec:swap}

The last example needs to change the caller's variables:
\begin{lstlisting}
(define r 100)
(define s 20)
(swap r s)        ; afterwards r should be 20 and s 100
\end{lstlisting}

\subsection{As a Function: It Does Nothing}

\begin{lstlisting}
(define (swap x y)
  (let ([t y])
    (set! y x)
    (set! x t)))
\end{lstlisting}
\code{(swap r s)} runs without error, but \code{r} and \code{s} keep their values. Racket passes the \emph{values} 100 and 20. The parameters \code{x} and \code{y} are new variables, with their own locations in the store in the sense of \Cref{ch:week6}, initialized with copies of those values. The two \code{set!}s update the parameters' locations; \code{r}'s and \code{s}'s locations are never touched.

\subsection{With Boxes: Pass the Locations}

If the caller hands over boxes instead of numbers, the function can mutate what the boxes contain:
\begin{lstlisting}
(define (swap-box b1 b2)
  (let ([t (unbox b2)])
    (set-box! b2 (unbox b1))
    (set-box! b1 t)))

(define r-box (box 10))
(define s-box (box 20))
\end{lstlisting}
After \code{(swap-box r-box s-box)}, \code{(unbox r-box)} is 20 and \code{(unbox s-box)} is 10. The parameters are still copies, but copies of box values, and a box value is a location: \code{b1} and \code{r-box} refer to the same box. This works, but only for data that the caller has put in boxes beforehand.

\subsection{As a Macro: Rewrite the Caller's Code}

\begin{lstlisting}
(define-syntax (swap stx)
  (syntax-case stx ()
    [(swap x y)
     #'(let ([t y])
         (set! y x)
         (set! x t))]))
\end{lstlisting}
Now \code{(swap r s)} expands, at the place where it is used, into
\begin{lstlisting}
(let ([t s])
  (set! s r)
  (set! r t))
\end{lstlisting}
The \code{set!}s name \code{r} and \code{s} themselves, so they update the caller's variables. No values are passed: the macro inserts code that refers to the variables directly.

\begin{keypoint}
A function receives values, so it can mutate only what it can reach from those values (the contents of a box, the fields of a mutable struct). A macro receives the caller's code, so it can expand into \code{set!} on the caller's own variables. This is the closest Racket comes to call-by-reference.
\end{keypoint}

\begin{remark}
The template introduces \code{t}. If the user's variable is also named \code{t}, as in \code{(swap t s)}, a text macro would expand to \code{(let ([t s]) (set! s t) (set! t t))}, which swaps nothing. Racket's hygienic expander keeps the two \code{t}s apart (\Cref{sec:hygiene}), and the swap works. This is the \code{SWAP} example of Recitation~5 (\Cref{sec:rec5}).
\end{remark}

\section{Breaking Hygiene on Purpose: \texttt{aif}}
\label{sec:aif}

This section was not covered in lecture. It comes from the end of the posted lecture code.

Hygiene (\Cref{sec:hygiene}) is almost always what we want, but occasionally a macro is \emph{meant} to bind a name that the user then writes. The standard example is the \emph{anaphoric if}. A common pattern is to compute a value, test it, and use it in a branch. With a division that returns \code{\#f} instead of failing,
\begin{lstlisting}
(define (safe-div x y)
  (if (= y 0) #f (/ x y)))

(let ([q (safe-div 10 2)])
  (if q (+ q 1) 'undefined))
\end{lstlisting}
the \code{let} is boilerplate. An anaphoric if, \code{(aif test then-e else-e)}, binds the value of \code{test} to the fixed name \code{it}, so that the branches can refer to it: \code{(aif (safe-div 10 2) (+ it 1) 'undefined)}.

\subsection{The Hygienic Version Fails}

The obvious definition is
\begin{lstlisting}
(define-syntax (aif-broken stx)
  (syntax-case stx ()
    [(aif-broken test then-e else-e)
     #'(let ([it test])
         (if it then-e else-e))]))
\end{lstlisting}
and \code{(aif-broken (safe-div 10 2) (+ it 1) 'undefined)} fails with \code{it: unbound identifier}. This is hygiene doing its job, in the first direction of \Cref{sec:hygiene}: the \code{it} bound in the template comes from the macro, the \code{it} in \code{(+ it 1)} comes from the user, and a binding introduced by the template does not capture a name the user wrote. The user's \code{it} is looked up where it was written, where nothing named \code{it} is bound.

\subsection{Choosing the Context of a Name}

To make the binding visible to the user's code, the macro has to create an \code{it} that counts as written by the user:
\begin{lstlisting}
(define-syntax (aif stx)
  (syntax-case stx ()
    [(aif test then-e else-e)
     (with-syntax ([it (datum->syntax stx 'it)])
       #'(let ([it test])
           (if it then-e else-e)))]))
\end{lstlisting}
\begin{itemize}
  \item \code{(datum->syntax stx 'it)} turns the symbol \code{'it} into a syntax object, taking its lexical context from \code{stx}, the use of \code{aif} that the user wrote. The new identifier therefore behaves as if the user had written \code{it} at the use site.
  \item \code{with-syntax} binds the pattern variable \code{it} to that syntax object, as if it had been matched by a pattern. In the template, both occurrences of \code{it} are pattern-variable references, so both are replaced by the user-context \code{it}, and the \code{let} now binds the same \code{it} that the branches refer to.
\end{itemize}
The transformer now runs code of its own before building the result, which is the case where \code{define-syntax-rule} is not enough and the full \code{syntax-case} form is needed (\Cref{sec:myor}). As intended,
\begin{lstlisting}
(aif (safe-div 10 2) (+ it 1) 'undefined)    ; 6
(aif (safe-div 10 0) (+ it 1) 'undefined)    ; 'undefined
\end{lstlisting}

\subsection{The Price}

The capture is now real, and it is not limited to the branches the macro author had in mind. It also hides any \code{it} the user already has:
\begin{lstlisting}
(let ([it "user's it"])
  (aif #f 3 it))
\end{lstlisting}
The \code{else} branch looks as if it uses the user's \code{it}, but it is inside the binding made by \code{aif}, so it gets the test value, and the result is \code{\#f}, not \code{"user's it"}. This is the capture problem of the C preprocessor (\Cref{sec:rec5}), now created on purpose.

\begin{keypoint}
In a hygienic system, capture is the exception and has to be asked for explicitly, one name at a time, with \code{datum->syntax}. A macro that does so makes that name part of its interface: users must know that \code{aif} binds \code{it}, just as they know that \code{lambda} binds its parameters.
\end{keypoint}

\section{Recitation 8: The Lambda Calculus}
\label{sec:rec8}

Our interpreters have kept growing: numbers, booleans, \code{let}, boxes, \code{rec}. This recitation goes the other way. It starts from a language with \emph{only} functions, the \emph{lambda calculus} (LC), and builds booleans, numbers, and recursion out of functions alone. The handout \code{recitation8.rkt} is written in \code{\#lang scribble/lp2}, a literate-programming language: the text and the code chunks render as a document, and the chunks run as a program with \code{check-expect} tests.

\subsection{Syntax}

LC has three kinds of terms:
\[
\mathit{LC} ::= x \mid \lambda x.\, e \mid e_1\ e_2
\]
written as s-expressions \code{(var x)}, \code{(lam x e)}, and \code{(app e1 e2)}. The identity function $\lambda x.\, x$ is \code{'(lam x (var x))}, and applying it to $y$ is \code{'(app (lam x (var x)) (var y))}.

Terms are built and taken apart with \emph{quasiquote}. The backquote quotes a list, and each comma plugs a Racket value into it:
\begin{lstlisting}
`(app ,e1 ,e2)         ; = (list 'app e1 e2)
\end{lstlisting}
Inside a \code{match} pattern, the same notation takes a list apart: \code{\textasciigrave(app ,e1 ,e2)} matches any three-element list whose first element is \code{'app}, and binds \code{e1} and \code{e2}. A pattern \code{,(? symbol?)} matches any symbol without binding it.

\subsection{Given: Renaming, Alpha-Equivalence, Substitution}

The handout provides three functions on terms, which the exercises build on.

\subsubsection{\code{is-lc?}} checks that an s-expression is a term, so \code{'(lam x x)} (body not a term) and \code{'(var 5)} (not a symbol) are rejected.

\subsubsection{\code{(rename-var lc x x-new)}} renames the \emph{free} occurrences of \code{x} to \code{x-new}. It stops at a \code{lam} that binds \code{x} again, since the occurrences under it refer to that inner binder:
\begin{lstlisting}
(define (rename-var lc x x-new)
  (match lc
    [`(var ,v) (if (eq? v x) `(var ,x-new) lc)]
    [`(lam ,v ,body)
     (if (eq? v x)
         lc
         `(lam ,v ,(rename-var body x x-new)))]
    [`(app ,e1 ,e2) `(app ,(rename-var e1 x x-new) ,(rename-var e2 x x-new))]))
\end{lstlisting}
\code{x-new} must be \emph{fresh}, that is, appear nowhere in \code{lc}. Otherwise it can be captured: renaming \code{x} to \code{y} in \code{'(lam y (var x))} would give \code{'(lam y (var y))}, a different function. The handout only calls \code{rename-var} with names from \code{gensym}, which are always fresh.

\subsubsection{Alpha-equivalence} \code{'(lam x (var x))} and \code{'(lam y (var y))} are the same function; they differ only in the name of the bound variable.

\begin{definition}[Alpha-equivalence]
Two terms are \emph{alpha-equivalent} if they are equal up to a consistent renaming of bound variables.
\end{definition}

\code{equiv-LC} checks this by renaming the binders of two \code{lam}s to one fresh name and comparing the bodies:
\begin{lstlisting}
(define (equiv-LC lc1 lc2)
  (match* (lc1 lc2)
    [(`(var ,x1) `(var ,x2)) (eq? x1 x2)]
    [(`(lam ,x1 ,body1) `(lam ,x2 ,body2))
     (let ([x-new (gensym)])
       (equiv-LC (rename-var body1 x1 x-new) (rename-var body2 x2 x-new)))]
    [(`(app ,e11 ,e12) `(app ,e21 ,e22))
     (and (equiv-LC e11 e21) (equiv-LC e12 e22))]
    [(_ _) #f]))
\end{lstlisting}
\code{match*} matches several values at once. Free variables must be literally equal: \code{'(lam x (var y))} is not equivalent to \code{'(lam y (var y))}, because the first body mentions a free \code{y} and the second the bound one. Nested binders may also swap names, so $\lambda x.\, \lambda y.\, x\ y$ and $\lambda y.\, \lambda x.\, y\ x$ are equivalent.

\subsubsection{Substitution} \code{(subst lc x arg)} replaces the free occurrences of \code{x} in \code{lc} by \code{arg}, avoiding capture:
\begin{lstlisting}
(define (subst lc x arg)
  (match lc
    [`(var ,v) (if (eq? v x) arg lc)]
    [`(lam ,v ,body)
     (if (eq? v x)
         lc
         (let ([v-new (gensym v)])
           `(lam ,v-new ,(subst (rename-var body v v-new) x arg))))]
    [`(app ,e1 ,e2) `(app ,(subst e1 x arg) ,(subst e2 x arg))]))
\end{lstlisting}
This is the capture-avoiding substitution of \Cref{sec:capavoid}, with a simplification: it renames the binder of \emph{every} \code{lam} it enters to a fresh name, whether or not capture could happen. Then no binder can capture a free variable of \code{arg}, and the code never has to compute which variables are free in \code{arg}. The cost is that results contain names like \code{y1825}, so tests compare them with \code{equiv-LC} instead of \code{equal?}. The capture case is
\begin{lstlisting}
(subst '(lam y (var x)) 'x '(var y))     ; alpha-equivalent to '(lam w (var y))
\end{lstlisting}
which must not become the identity function \code{'(lam y (var y))}.

\subsection{Evaluation: Normal Order}

\begin{definition}[Redex, normal form]
A \emph{redex} is a function applied to an argument, \code{\textasciigrave(app (lam ,x ,body) ,arg)}. \emph{Reducing} it means substituting \code{arg} for \code{x} in \code{body}. A term with no redex anywhere in it is in \emph{normal form}, and the normal form is the answer.
\end{definition}

The interpreter of this recitation evaluates in \emph{normal order}, which differs from our earlier interpreters in three ways:
\begin{itemize}
  \item It reduces the \emph{outermost} redex first. \code{(app (lam x (var z)) (app (lam y (var y)) (var b)))} reduces at once to \code{(var z)}, without touching the argument. Racket and our interpreters would evaluate the argument first.
  \item Among subterms, it goes \emph{left to right}.
  \item It also reduces \emph{under} \code{lam}. Our interpreters treated a function as a value and never looked inside it; here \code{(lam x (app (lam y (var y)) (var x)))} steps to \code{(lam x (var x))}.
\end{itemize}

\begin{exercise}[One step]
Write \code{(step-LC lc)}, which finds the redex that normal order picks, reduces it, and returns the new term, or returns \code{\#f} if \code{lc} has no redex. Before writing code, fill in the expected results of the handout's tests:
\begin{lstlisting}
(step-LC '(app (lam x (var x)) (var y)))
(step-LC '(var x))
(step-LC '(app (var f) (var a)))
(step-LC '(app (lam x (var z)) (app (lam y (var y)) (var b))))
(step-LC '(app (app (lam x (var x)) (var a)) (app (lam y (var y)) (var b))))
(step-LC '(lam x (app (lam y (var y)) (var x))))
\end{lstlisting}
Hints from the handout: for an \code{app} that is not a redex, try to step \code{e1}, and if that returns \code{\#f}, try \code{e2}. The redex case must come before the general \code{app} case in the \code{match}, since every redex also matches \code{\textasciigrave(app ,e1 ,e2)}.
\end{exercise}
\begin{solution}
The expected results, in order:
\begin{lstlisting}
'(var y)                                       ; the redex itself
#f                                             ; a variable has no redex
#f                                             ; f is not a lam
'(var z)                                       ; outermost first; argument dropped
'(app (var a) (app (lam y (var y)) (var b)))   ; left subterm first
'(lam x (var x))                               ; under the lam
\end{lstlisting}
The fifth term is not itself a redex, since its function part is an \code{app}, not a \code{lam}. So the left subterm steps, and the right one is left for later. In the last test, \code{subst} enters no \code{lam}, so the result has no gensym names and \code{equal?} works.

\code{step-LC} has one case per form, with the redex first:
\begin{lstlisting}
(define (step-LC lc)
  (match lc
    [`(app (lam ,x ,body) ,arg) (subst body x arg)]
    [`(app ,e1 ,e2)
     (let ([s1 (step-LC e1)])
       (if s1
           `(app ,s1 ,e2)
           (let ([s2 (step-LC e2)])
             (and s2 `(app ,e1 ,s2)))))]
    [`(lam ,x ,body)
     (let ([s (step-LC body)])
       (and s `(lam ,x ,s)))]
    [`(var ,_) #f]))
\end{lstlisting}
\code{(and s2 \dots)} is \code{\#f} when \code{s2} is \code{\#f}, and the rebuilt term otherwise. Each case encodes one rule of normal order. The redex case reduces the outermost redex before looking inside the argument. The \code{app} case goes left to right, trying \code{e2} only when \code{e1} has no redex. The \code{lam} case steps under the binder.
\end{solution}

\begin{exercise}[Normal form]
Write \code{(interp-LC lc)}, which steps until no redex is left and returns the normal form.
\end{exercise}
\begin{solution}
\begin{lstlisting}
(define (interp-LC lc)
  (let ([next (step-LC lc)])
    (if next
        (interp-LC next)
        lc)))
\end{lstlisting}
\code{step-LC} returns \code{\#f} exactly when \code{lc} is in normal form, and that is the answer. The call is in tail position, so a long reduction does not grow the stack.
\end{solution}

Not every term has a normal form. The classic example is
\[
\Omega = (\lambda w.\, w\ w)\ (\lambda w.\, w\ w),
\]
\code{omega} in the handout. Its only redex substitutes the second copy for \code{w} in \code{(app (var w) (var w))}, which gives back $\Omega$ itself, so \code{(interp-LC omega)} runs forever. The test \code{(check-expect (step-LC omega) omega)} can use \code{equal?}: the body \code{(app (var w) (var w))} contains no \code{lam}, so \code{subst} renames nothing and the result is literally \code{omega} again.

The handout's tests for \code{interp-LC} include
\begin{lstlisting}
(interp-LC `(app (lam x (var z)) ,omega))      ; should be '(var z)
\end{lstlisting}

\begin{exercise}[Discussion]
What would an interpreter that evaluates arguments first, like Racket, do with \code{\textasciigrave(app (lam x (var z)) ,omega)}?
\end{exercise}
\begin{solution}
It runs forever. Before calling \code{(lam x (var z))}, it evaluates the argument $\Omega$, which steps only to itself. Normal order calls the function first. The function never uses \code{x}, so $\Omega$ is thrown away without being evaluated. In general, if a term has a normal form, normal order is guaranteed to find it (the \emph{standardization theorem}). Evaluating arguments first can loop on such a term, as here.
\end{solution}

\subsection{Church Booleans}

LC has no booleans, but they can be \emph{encoded} as functions. A boolean is a function that chooses between two arguments, \emph{true} picking the first and \emph{false} the second:
\[
\mathit{true} = \lambda x.\, \lambda y.\, x \qquad \mathit{false} = \lambda x.\, \lambda y.\, y
\]
\begin{lstlisting}
(define LC-true  '(lam x (lam y (var x))))
(define LC-false '(lam x (lam y (var y))))
\end{lstlisting}
\code{NF-to-bool} reads a normal form back as a Racket boolean, comparing with \code{equiv-LC}, and signals an error for anything else.

\begin{exercise}[If]
Define \code{LC-if}, a term that takes a boolean \code{c} and two terms \code{t} and \code{e}, one argument at a time, and returns \code{t} if \code{c} is true and \code{e} otherwise. The tests apply it as \code{\textasciigrave(app (app (app ,LC-if ,LC-true) (var yes)) (var no))}.

Hint: a boolean already \emph{is} a choice between two arguments. What is left for \code{LC-if} to do?
\end{exercise}
\begin{solution}
Nothing but passing the branches to the boolean: $\mathit{if} = \lambda c.\, \lambda t.\, \lambda e.\, c\ t\ e$.
\begin{lstlisting}
(define LC-if
  '(lam c (lam t (lam e (app (app (var c) (var t)) (var e))))))
\end{lstlisting}
With $c = \mathit{true}$, $c\ t\ e$ reduces to $t$. Under normal order the branch not taken is never evaluated, so \code{LC-if} needs no special treatment. In Racket it does: there, \code{if} must be a special form, and a function \code{my-if} fails (\Cref{sec:myor}).
\end{solution}

\begin{exercise}[Not, And]
Define \code{LC-not} and \code{LC-and} in the same style. \code{LC-true} and \code{LC-false} can be plugged in with quasiquote.

Hint: write each one as ``if \dots{} then \dots{} else \dots'' first, and then use that a boolean is its own \code{if}.
\end{exercise}
\begin{solution}
$\mathit{not}\ b$ is ``if $b$ then false else true'', and $\mathit{and}\ a\ b$ is ``if $a$ then $b$ else false''. Since a boolean applied to two terms is already an \code{if}:
\begin{lstlisting}
(define LC-not
  `(lam b (app (app (var b) ,LC-false) ,LC-true)))
(define LC-and
  `(lam a (lam b (app (app (var a) (var b)) ,LC-false))))
\end{lstlisting}
The commas plug the terms \code{LC-false} and \code{LC-true} in. Without the quasiquote, \code{'(\dots{} LC-false \dots)} would contain the symbol \code{LC-false}, which is not a term. \code{LC-and} returns $b$ itself when $a$ is true, just as \code{and} in Racket returns its last argument.
\end{solution}

\subsection{Church Numerals}

The natural number $n$ is encoded as the function that applies $f$ to $x$, $n$ times:
\[
\bar 0 = \lambda f.\, \lambda x.\, x \qquad
\bar 1 = \lambda f.\, \lambda x.\, f\ x \qquad
\bar 2 = \lambda f.\, \lambda x.\, f\ (f\ x) \qquad
\bar 3 = \lambda f.\, \lambda x.\, f\ (f\ (f\ x))
\]
These are \emph{Church numerals}. \code{(LC-nat n)} builds the numeral for a Racket number, and \code{NF-to-nat} reads a normal form back by counting the \code{f}s. It works for any names of the two binders, so the gensym names in a result do not matter.

\begin{exercise}[Arithmetic]
Define \code{LC-succ}, \code{LC-add}, and \code{LC-mul}. Each takes numerals and returns a numeral, a term of the form $\lambda f.\, \lambda x.\, \dots$. Hints from the handout:
\begin{itemize}
  \item $n\ f\ x$ applies $f$ to $x$, $n$ times. Successor applies it once more.
  \item $m + n$: apply $f$ $n$ times, then $m$ more times.
  \item $m \times n$: $n\ f$ is a function that applies $f$ $n$ times. Apply \emph{that} $m$ times.
\end{itemize}
The tests include $0 \times 7$ and $3 \times 0$.
\end{exercise}
\begin{solution}
Each definition takes its numerals, then $f$ and $x$, and builds the body from the hint:
\begin{align*}
\mathit{succ} &= \lambda n.\, \lambda f.\, \lambda x.\, f\ (n\ f\ x) \\
\mathit{add} &= \lambda m.\, \lambda n.\, \lambda f.\, \lambda x.\, m\ f\ (n\ f\ x) \\
\mathit{mul} &= \lambda m.\, \lambda n.\, \lambda f.\, \lambda x.\, m\ (n\ f)\ x
\end{align*}
\begin{lstlisting}
(define LC-succ
  '(lam n (lam f (lam x (app (var f) (app (app (var n) (var f)) (var x)))))))
(define LC-add
  '(lam m (lam n (lam f (lam x (app (app (var m) (var f))
                                    (app (app (var n) (var f)) (var x))))))))
(define LC-mul
  '(lam m (lam n (lam f (lam x (app (app (var m) (app (var n) (var f)))
                                    (var x)))))))
\end{lstlisting}
In \code{LC-add}, $n\ f\ x$ is $f$ applied $n$ times to $x$, and $m\ f$ applies $f$ $m$ more times to that. In \code{LC-mul}, $n\ f$ is ``apply $f$ $n$ times'', and $m$ repeats that $m$ times, $m \cdot n$ applications in total. The zero cases need no special handling. For $0 \times 7$, $\bar 0$ ignores its function and returns $x$. For $3 \times 0$, $\bar 0\ f$ is the identity, and three identities applied to $x$ give $x$.

Because normal order reduces under \code{lam}, the result is a numeral in normal form, which is what \code{NF-to-nat} expects. The shorter $\mathit{mul} = \lambda m.\, \lambda n.\, \lambda f.\, m\ (n\ f)$ also passes the tests: $m$ is a numeral, so $m\ (n\ f)$ reduces to a $\lambda x.\, \dots$ itself.
\end{solution}

\subsubsection{Given: zero test and predecessor} \code{LC-is-zero} applies the constant function $\lambda u.\, \mathit{false}$ to $\mathit{true}$, $n$ times. Zero applications leave $\mathit{true}$, and one or more give $\mathit{false}$:
\begin{lstlisting}
(define LC-is-zero
  `(lam n (app (app (var n) (lam u ,LC-false)) ,LC-true)))
\end{lstlisting}

\code{LC-pred} is famously tricky, and it is given:
\begin{lstlisting}
(define LC-pred
  '(lam n
     (lam f
       (lam x
         (app (app (app (var n)
                        (lam g (lam h (app (var h)
                                           (app (var g) (var f))))))
                   (lam u (var x)))
              (lam u (var u)))))))
\end{lstlisting}
It builds the result one $f$ at a time, but its starting value $\lambda u.\, x$ throws away the first $f$, so $n$ copies of $f$ become $n - 1$. The predecessor of $\bar 0$ is $\bar 0$. Wikipedia's ``Church encoding'' page has the details.

\subsection{Challenge: Recursion with \texttt{Y}}

In Recitation~7 (\Cref{sec:rec7}), \code{Y} looped forever in Racket and \code{Z} was needed instead. This part tries \code{Y} again, in LC with normal order.

The handout provides Racket helpers that build bigger terms, so that not every \code{app} has to be written by hand. Their arguments are LC terms, so write \code{(SLC-zero? '(var n))}, not \code{(SLC-zero? 'n)}:
\begin{lstlisting}
(define (SLC-if c t e)  `(app (app (app ,LC-if ,c) ,t) ,e))
(define (SLC-zero? e)   `(app ,LC-is-zero ,e))
(define (SLC-pred e)    `(app ,LC-pred ,e))
(define (SLC-mul e1 e2) `(app (app ,LC-mul ,e1) ,e2))
\end{lstlisting}
\code{SLC-if} and \code{SLC-mul} use your \code{LC-if} and \code{LC-mul}, so if the factorial tests fail, check those first.

\begin{exercise}[Y and factorial]
Write \code{LC-Y}, the Y combinator of \Cref{sec:rec7},
\[
Y = \lambda h.\, (\lambda x.\, h\ (x\ x))\ (\lambda x.\, h\ (x\ x)),
\]
as an LC term. Then write \code{LC-fact-gen}, the LC version of \code{fact-gen}: it takes the recursive call \code{fact} and returns a function of \code{n}. The tests check $0! = 1$, $3! = 6$, and $4! = 24$ via \code{\textasciigrave(app (app ,LC-Y ,LC-fact-gen) ,(LC-nat 3))}.
\end{exercise}
\begin{solution}
\code{LC-Y} is the formula transcribed:
\begin{lstlisting}
(define LC-Y
  '(lam h (app (lam x (app (var h) (app (var x) (var x))))
               (lam x (app (var h) (app (var x) (var x)))))))
\end{lstlisting}
\code{LC-fact-gen} is \code{fact-gen} of \Cref{sec:rec7}, built with the helpers:
\begin{lstlisting}
(define LC-fact-gen
  `(lam fact
     (lam n
       ,(SLC-if (SLC-zero? '(var n))
                (LC-nat 1)
                (SLC-mul '(var n)
                         `(app (var fact) ,(SLC-pred '(var n))))))))
\end{lstlisting}
The comma before \code{(SLC-if \dots)} switches back to Racket, so the helpers run and build the body of the \code{lam}. The inner backquote starts a new quasiquote for the recursive call \code{(fact (pred n))}. All three tests pass. In the Python port used to check these solutions, $4!$ takes about 4100 steps.
\end{solution}

\begin{exercise}[Discussion]
\leavevmode
\begin{enumerate}
  \item Why does \code{Y} work here, when it looped in Racket last week?
  \item Does \code{(interp-LC \textasciigrave(app ,LC-Y ,LC-fact-gen))}, without a number, finish?
  \item Write \code{Z} as an LC term. Does it compute factorial in LC? What about \code{Z} applied to \code{LC-fact-gen} alone, without a number?
\end{enumerate}
\end{exercise}
\begin{solution}
\begin{enumerate}
  \item In Racket, \code{(Y fact-gen)} loops because the argument \code{(W W)} must be evaluated before \code{fact-gen} is called, and \code{(W W)} only unfolds into itself (\Cref{sec:rec7}). Normal order reduces the outermost redex first, so \code{fact-gen} is called with \code{(W W)} \emph{unevaluated}, as in the lazy evaluation of \Cref{sec:lazy}. \code{(W W)} unfolds one more level only when it reaches function position, that is, when factorial actually makes a recursive call. At $n = 0$, the \code{if} picks the base case, and the remaining copy of \code{(W W)} is discarded.
  \item No. $Y\ \mathit{fact\mbox{-}gen}$ reduces to $\lambda n.\, \mathit{if}\ (\mathit{iszero}\ n)\ \bar 1\ (n \times (W\ W)\ (\mathit{pred}\ n))$. Here $n$ is a variable, so the test cannot decide between the branches. But normal order reduces under $\lambda$, and so inside both branches. In the second branch, $(W\ W)$ unfolds into another copy of the factorial body, which contains another $(W\ W)$, and so on. The term has no normal form, and \code{interp-LC} runs forever. With a number, the test decides, and the unfolding stops.
  \item \code{Z} is \code{Y} with each \code{(x x)} wrapped in \code{(lam v \dots)}:
\begin{lstlisting}
(define LC-Z
  '(lam h (app (lam x (app (var h) (lam v (app (app (var x) (var x)) (var v)))))
               (lam x (app (var h) (lam v (app (app (var x) (var x)) (var v))))))))
\end{lstlisting}
  With a number, it computes factorial as well: $0!$ and $3!$ come out right, and $3!$ takes about the same number of steps as with \code{Y}. Without a number, it also runs forever. The \code{(lam v \dots)} delayed \code{(x x)} in Racket because Racket never evaluates inside a \code{lambda}. Normal order does reduce under \code{lam}, so the wrapper delays nothing, and \code{(x x)} unfolds forever as before.
\end{enumerate}
The lesson: \code{Y} versus \code{Z} is not a property of the combinators alone. It depends on the evaluation strategy, namely whether arguments are evaluated before the call and whether evaluation goes under $\lambda$.
\end{solution}

